\section{Discussion}\label{sec:discussion}
\subsection{Evaluation Required for Predictive Claims}\label{sec:future}
The next empirical task is to define one measurable outcome, resolution source,
and horizon before issuing forecasts. For example, a creator-cash task must
specify withdrawals, expenses, token inventory, and liabilities. Until those
quantities are joined to the strategic model and observed outcomes, a claimed
probability of creator profit is not a supported output.

A prospective study records predictions before outcomes are known. It must
freeze the construction procedure, evidence cutoff, rules, model version, and
selection mode for each evaluation release. Training, development, and holdout
sets must follow time order, and overlapping market events need separation to
avoid leaking outcome information. A holdout set is data reserved for the final
assessment rather than prompt or parameter tuning. Failed runs and abstentions
must remain in the issued-prediction record.

Comparisons should include a training-period base rate, a simple numeric rule,
and a single-model forecast given the same available evidence. Controlled economic-game tasks should first compare choice behavior with known
payoffs and supplied information, following the motivation of OpenSpiel and
repeated-game studies \cite{lanctot2019,akata2025}. A matched-budget
comparison tests whether interacting participants add value beyond spending
more inference resources. Ablations should remove peer messages, replace
heterogeneous goals with identical roles, and compare Jev with a local utility
rule and a larger-model policy. Additional participants are not independent
observations when they share assumptions or correlated model errors. Prompt
experiments should compare a minimal typed contract with numeric-support
criteria, vary menu order, and test whether selected actions respond to
controlled payoff and information changes. Tuning must use development cases,
not future resolved outcomes or the final holdout set. Quality and latency
must be measured alongside actual billed input use; neither fluent motives
nor lower per-call prices establish useful strategic behavior.

\begin{figure*}[t]
\centering
\begin{tikzpicture}[node distance=6mm and 8mm]
\node[box] (task) {Define outcome and horizon\\Fix evidence cutoff};
\node[box,right=of task] (freeze) {Freeze evaluation protocol\\Rules, models, baselines};
\node[box,right=of freeze] (issue) {Issue dated forecasts\\Preserve abstentions};
\node[box,below=of issue] (resolve) {Resolve observed outcomes\\Use the declared source};
\node[box,left=of resolve] (score) {Compare held-out scores\\Cost, coverage, uncertainty};
\node[box,left=of score] (publish) {Report results and failures\\Limit claims to evidence};
\draw[arr] (task)--(freeze);\draw[arr] (freeze)--(issue);
\draw[arr] (issue)--(resolve);\draw[arr] (resolve)--(score);\draw[arr] (score)--(publish);
\end{tikzpicture}
\caption{Proposed prospective evaluation. This is a research protocol for
future validation, not an experiment completed in this paper.}
\label{fig:evaluation}
\end{figure*}

Evaluation should report Brier score, log loss, calibration, completion rate,
latency, and actual token usage. Calibration asks whether events assigned a
given probability occur at approximately that frequency across sufficiently
many resolved cases. Cases sharing a narrative or time period must be treated
as dependent when estimating uncertainty. Repeated samples from one assumed
world do not substitute for independent real outcomes.

Decision usefulness requires a separate user study. Invited teams could be
asked to formulate a hypothesis, identify a consequential assumption,
distinguish utility from cash,
and explain whether the report changes their decision. Completion time and
consented feedback can assess practical value. This study tests deliberation,
while resolved-outcome scoring tests prediction; one result cannot establish
the other.

\subsection{Research Applications and Operating Cost}\label{sec:business}
The initial users are researchers, protocol designers, and risk analysts
studying financial and economic decisions, initially bounded crypto and DeFi
questions. Each included actor must have an explicit goal, feasible options,
resources and incentives that matter to the mechanism being studied. Typical tasks compare an
incentive mechanism, change an information policy, or test a stress assumption
while holding the rest of an experiment fixed. The value hypothesis is a
shorter path from a research question to an inspectable experiment, supported
by versioned assumptions, sensitivity comparisons, and reusable results.
Spreadsheets, scripted agent-based models, and single-model analysis are useful
baselines. Demand and willingness to pay remain questions for the user study.

Typed small-model inference offers an economical way to evaluate bounded
choices, but the cost example in Section~\ref{sec:system} covers only actor
input tokens. World authoring, hosting, storage, failed requests, and support
must be measured separately. No matched-task comparison currently supports a
claim of lower cost at equal quality than MiroFish or a frontier model.
The relevant operating measures are cost per completed, replayable experiment
and cost per evaluated research hypothesis. Low per-call prices alone do not
establish a useful research tool. Provider interface documentation also does
not establish commercial resale rights or a particular privacy guarantee.

\subsection{Remote Agent Services}
The local MCP interface enables an external agent to construct a world with
its own model and invoke a bounded engine. Distributing those services would
require additional identity, authorization, and reliability mechanisms.
Proof of Response proposes responses, fault evidence, or delay compensation
across a network of service dependencies \cite{proofresponse2025}. Boyun's
recorded timeout and cooldown handling does not implement that protocol.
Similarly, a response delivered on time can still contain an invalid decision
or a poor prediction. Service responsiveness and scientific validity require
separate evaluation.

\subsection{Limitations}\label{sec:limits}
\emph{Behavioral assumptions.} Motivations, utility weights, action effects,
and population composition are predominantly reviewed hypotheses. Wallet
observations are incomplete and cannot identify private intent or account
ownership. Detailed personas can share the same model bias. Their number
does not establish a representative market population.

\emph{Strategic simplification.} Linear utility, additive metric effects,
scheduled events, and preset publications omit many nonlinearities and
adaptive behaviors. Repeated rounds do not imply repeated-game equilibrium,
and public claims are not enforceable commitments. The strategic engine does
not model a full exchange, credit, leverage, or liquidation process.

\emph{Information and privacy.} Typed responses prevent malformed action
execution but do not eliminate mistaken reasoning or prompt injection, where
untrusted text attempts to redirect the model. Restricted observations can
still reveal information through public consequences. The host can read the
world and remote providers receive their selected requests. Artifact
commitments establish consistency, not anonymity or confidential computation.

\emph{External validity.} The synthetic experiments verify specified
calculations and scheduling. The live request verifies one interface path.
Neither establishes human fidelity, prospective forecast accuracy, or an
improvement in real decisions. Those claims require the comparative and
prospective studies described above.

\section{Conclusion}\label{sec:conclusion}
Boyun separates authored assumptions, typed participant choices, and numeric
consequences in a replayable simulation workflow for financial and economic decisions. Its world specification
makes individual goals, incentives, information boundaries, resource costs, and simultaneous
interaction rules explicit. Six synthetic experiments verify arithmetic,
timing, and mechanism attribution under declared test policies. Optional NEAR
registration supports later artifact-integrity checking while keeping the
artifact contents off chain.

The remaining research question is whether this structure improves predictions
or decisions relative to simpler alternatives. A frozen prospective evaluation
with baselines and ablations, together with a separate user study, is needed
to answer it. Those results should determine which assumptions and interaction
mechanisms are useful enough to retain.
